\documentclass[11pt,a5paper]{article}
\usepackage[margin=17mm,headsep=5mm]{geometry}
\usepackage{amsmath,amssymb,mathtools,xcolor,fancyhdr,booktabs,tabularx,titlesec}
\usepackage[colorlinks=true,linkcolor=blue!46!black,urlcolor=blue!45!black]{hyperref}

\usepackage{fontspec}
\setmainfont[AutoFakeSlant=0.16]{Noto Serif CJK SC}
\setsansfont{Noto Sans CJK SC}
\XeTeXlinebreaklocale "zh"
\XeTeXlinebreakskip=0pt plus 0.15em
\renewcommand{\contentsname}{阅读路线}
\setlength{\parindent}{2em}
\setlength{\parskip}{.28em}
\linespread{1.20}
\allowdisplaybreaks[2]\setlength{\emergencystretch}{1em}
\definecolor{ink}{HTML}{17344B}
\titleformat{\section}{\Large\bfseries\sffamily\color{ink}}{\thesection}{0.6em}{}\titleformat{\subsection}{\large\bfseries\sffamily\color{ink}}{\thesubsection}{0.6em}{}\pagestyle{fancy}\fancyhf{}
\fancyhead[L]{\small 数学证明审读}\fancyhead[R]{\small 从整除到黎曼猜想}
\fancyfoot[C]{\small\thepage}
\setlength{\headheight}{15pt}
\newcommand{\rk}{\operatorname{rk}}
\newcommand{\tr}{\operatorname{tr}}
\newcommand{\Gal}{\operatorname{Gal}}
\newcommand{\Span}{\operatorname{span}}
\newcommand{\src}[3]{\href{https://github.com/leanprover-community/riemann-hypothesis/blob/40e56bab026552e089d7afcfeff08d2defb9c005/src/#1\##2}{#3}}
\newcommand{\Rr}{\mathbb R}
\newcommand{\Cc}{\mathbb C}
\newcommand{\Zz}{\mathbb Z}
\begin{document}
\thispagestyle{empty}
{\sffamily\color{ink}\LARGE\bfseries 从整除到黎曼猜想\par}
\smallskip
{\sffamily\small 试读修订版 · 2026 年 11 月 26 日 · 第二次修订\par}
\medskip

\noindent\textbf{要证明什么？} 满足 $\zeta(s)=0$ 的 $s$ 叫零点。负偶数上的零点早已知道；黎曼猜想说，其余零点都在实部为 $1/2$ 的直线上。若能对每个 $\varepsilon>0$ 排除 $\operatorname{Re}s>1/2+\varepsilon$ 的非平凡零点，就排除了整条直线右侧：一个零点如果在右边，取它到直线距离的一半作 $\varepsilon$ 即可。再由零点关于临界线的对称性排除左侧。要任意缩小的是允许的\textbf{横向距离}，不是零点的高度。

\noindent\textbf{整条证明怎么走？} 先用一串 theta 提升把有关素数的和联系到一个格。这里的提升是把一类自守形式转到另一类的变换。格的余体积衡量一组独立向量张成的基本区域有多大；指数则衡量这组向量生成的子格漏了多少陪集。结合短向量、足够多的独立方向和整性给出的下界，两边估计同一个带指数修正的量。共有的主项抵消后，剩下对素数计数偏差的单边上界。再用解析判据排除右半平面的零点，由函数方程排除左侧。

\noindent\textbf{本章做哪一步？} 我们要讲清\textbf{整除为什么能限制这个偏差}：局部整除限制哪些向量能进入格，矩阵估计控制这些向量的大小和独立方向。只有两边确实针对同一个量，才可以把界合起来。

\noindent\textbf{先看一个整数。} 若 $A\ne0$，且 $125$ 和 $49$ 都整除 $A$，因为二者互素，$125\cdot49=6125$ 也整除 $A$，所以 $|A|\ge6125$。若又能证明同一个 $A$ 的绝对值小于 $6125$，两条信息就不能同时成立。整除给下限，另一边给上限；\textbf{比较的必须是同一个数}。

\vfill
{\small 本章解释这段推导，并标出所用的外部结果；不在这里独立重证整个黎曼猜想。第 8 节核对「所有输入都判真」的说法。正文沿用原版源码定位；下一页目录后附三版对照。\par}
\newpage
\setcounter{tocdepth}{1}
\hypersetup{bookmarksdepth=2}
\tableofcontents
\vspace{0.65cm}
\noindent\textbf{第一次读：} 先读第 1 节，找出要估计的偏差以及两条证据在哪里会合。第 2、3 节拆开讲整除、矩阵和斜率；第 5 节把它们合起来。第 4 节供想追查实际矩阵和源码的读者回看。

\medskip
\noindent\textbf{不必先学全套工具。} 开头只用整除、坐标和例子。格、对偶和斜率在用到时解释；解析公式单列，不是继续读代数部分的门槛。
\newpage
\section*{证明路线与三个源码版本}
\phantomsection
\addcontentsline{toc}{section}{证明路线与三个源码版本}

\noindent\textbf{从哪里读起？} 先看精简版的\href{https://github.com/lsernau/grh-lean/blob/72fead102f663d5c7a407964a13dc0cf9d977207/grh-compact/src/New/Main.lean\#L26-L72}{主证明文件开头}。把带指数修正的余体积估计记作 $H$，把双方共有的主项记作 $B$，把素数和偏差记作 $R$。暂时不写误差，比较的骨架是
\[
 B \le H \le B-mR,
\]
其中 $m$ 是提升次数。左边限制 $H$ 不能太小，右边说明 $R$ 若太大，$H$ 就会太小。消掉双方共有的 $B$，再带回误差，就得到 $R$ 的上界。选取第二个尺度并使用解析判据，才从这个界走到零点结论。

\noindent\textbf{原版。} Alpöge 与 Claude 发布的形式化工作。本文正文仍按\href{https://github.com/leanprover-community/riemann-hypothesis/tree/40e56bab026552e089d7afcfeff08d2defb9c005}{原版提交 \texttt{40e56bab}}定位。第 2.4 节另有 Chey 的固定链接；阶为 125 的手算例子只说明局部保阶，不是满足 Type I 筛选条件的原始算例。

\noindent\textbf{可读版。} Sernau 等人将原论证重写为约 30 万行代码。社区维护仓库的新版 \texttt{main} 从这个版本发展，适合按\href{https://github.com/leanprover-community/riemann-hypothesis/blob/2db90859500407542617a5b86281b4f1e1ee69cb/GUIDE.md}{社区阅读路线}定位定义和调用。本文引用原版时，文件名不能照搬到社区版。

\noindent\textbf{精简版。} 同一仓库的\href{https://github.com/lsernau/grh-lean/blob/72fead102f663d5c7a407964a13dc0cf9d977207/grh-compact/README.md}{\texttt{grh-compact}} 约 7.6 万行代码。它保留相同的五条顶层定理和上述证明骨架，但有些中间步骤换了论证。不能把内部文件逐项对号；两版都用 A、B、C 组织全局论证，但局部的 Hecke 整除估计和第二尺度的选取都换了做法。

\noindent\textbf{阅读建议。} 用精简版建立全局图，再回到可读版追查具体估计。本文各节仍会说清所依据的是哪一版；旧版中的例子不能直接当成精简版的实现。
\clearpage

\section{先认清要估计的数}
\subsection{用素数幂算一个偏差}
取一个尺度 $x$。对每个不超过 $x$ 的素数幂 $p^r$，按下面的规则记一项：
\[
       \left(\frac{x}{p^r}-1\right)\log p.
\]
素数幂越靠近 $x$，这项的权重越小；恰等于 $x$ 时贡献为零。比如 $x=6$，只用 $2,3,4,5$ 四个素数幂，算出来是
\[
 S(6)=2\log2+\log3+\tfrac12\log2+\tfrac15\log5.
\]
$4=2^2$ 那一项仍乘 $\log 2$，不是 $\log4$。把这条规则写成通式，令 $\Lambda(p^r)=\log p$，其余 $\Lambda(n)=0$，则
\[
 \begin{aligned}
 S(x)&=\sum_{1\le n\le x}\Lambda(n)(x/n-1),\\
 G(x)&=S(x)-\bigl[x\log x-(1+\gamma_E)x\bigr].
 \end{aligned}
\]
$\gamma_E$ 是欧拉常数。方括号是用作比较的主项；$G(x)$ 才是我们说的“偏差”，即实际算出的数减去主项。选择这种权重和主项，与 $\zeta$ 函数的积分公式有关，见本节末。\src{Development/AutomorphicArithmetic/DedekindLFunctions/CharFamily.lean}{L10-L22}{权重与偏差定义}。

本章不是直接证明对每个 $x$ 都有某个 $|G(x)|$ 上界。实际论证还要引入加权的偏差组合，并在合适的长区间内控制它的单侧上界。第 5 节写出区间条件；从它到零点还要用解析判据。先记住这里的分工：$G(x)$ 是解释最终对象的起点，后文的局部收益不是把 $G(x)$ 换个名字。

\subsection{两条证据在哪里用到一起}
先用整数矩阵把接口说清。假如一个矩阵的两列是 $(2,0)$、$(4,0)$，每个坐标都是整数，但两列只有一个独立方向。另一个矩阵的两列是 $(1/2,0)$、$(0,1/2)$，方向有两个，却不能把所有整数坐标送成整数坐标。前一个例子表明\textbf{整性不推出满秩}，后一个表明\textbf{满秩不推出整性}。

在真正的证明里，也把同一个映射写成矩阵。一条路证明：先定下同一个检验格，再按每个有限处的整除预算调整输入格，证明映射把调整后的输入送进指定的整对偶格。另一条路证明：用复矩阵控制这些映射的大小；删去一些输入方向后，把各个映射剩余的像合起来，仍能给出独立方向数的下界。“像”就是所有输出组成的空间，“秩”是其中独立方向的个数。

这两条路在第 3 节会合。那里要问：删掉一部分输入方向之后，剩下的像在\textbf{同一个输出格}里还占多大？要回答它，既要知道像有多少独立方向，又要知道这些方向确实来自指定格里的输出。换成另一个矩阵，即使各自的统计量都好看，也不能一起用。到了第 5 节，两个不等式从上下夹住同一个量：按各输入格的秩加权后的斜率。比较上下界时，同一个输出格的平均斜率消去，留下要控制的算术差。这就是矩阵那一支究竟补了什么。

\subsection{进入正文前，只认三个对象}
最简单的格是平面上的整数点 $\mathbb Z^2$。这里的格也规定哪些系数可以取，只是整数换成了数域的整数环。下面用到的三个对象是输入格、检验格和映射。做比较时固定一组数据。用 $u$ 编号不同的算术链，源码叫它“词”；这里先把它当作输入格的编号。

写 $R$ 为公共检验格，$C=R^\vee$ 为整对偶格。对偶这里的含义很具体：线性泛函取 $R$ 的任意元素，结果都在整数环。源侧从原格 $N_u^0$ 变成理想平移后的 $N_u$。只有在平移后才能写出本章要用的包含
\[
       B_u(N_u)\subseteq C.
\]
记 $h_u$ 为同一个 $B_u$ 的高度上界。证明不会从整性直接推出高度，也不会从高度推回整性。本章使用 $B_u$ 同时表示双线性配对与把第二输入留空后的线性映射，亦即 $\phi\mapsto (F\mapsto B_u(\phi,F))$。


\subsection{补读：为什么选择这一种权重}
熟悉复分析的读者可以用下式核对权重的来源；跳过本段不影响第 2 节。令 $W(t)=(t-1)_+$，在 $\Re z>1$ 时
\[
 \int_1^\infty(t-1)t^{-z-1}dt
  =\frac1{z-1}-\frac1z=\frac1{z(z-1)}.
\]
取 $c>1$。反演这个 Mellin 变换，再用等式
\[
 -\zeta'(z)/\zeta(z)=\sum_n\Lambda(n)n^{-z},
\]
便得到
\[
 S(x)=\frac{1}{2\pi i}\int_{c-i\infty}^{c+i\infty}
     -\frac{\zeta'(z)}{\zeta(z)}
       \frac{x^z}{z(z-1)}\,dz .
\]
这里可交换绝对收敛的级数与积分，因为在直线 $\Re z=c$ 上，两项条件
$\sum_n\Lambda(n)n^{-c}<\infty$ 与
$\int_\Rr |(c+it)(c-1+it)|^{-1}dt<\infty$ 同时成立。若 $x$ 恰是整数，边界权重为零，无需决定是否给最后一项半权。

这给出了一点具体的预期：$z=1$ 处，对数导数的极点与分母的极点重叠，二阶极点带出 $x\log x$ 型主项；若考察非平凡零点 $\rho$，其指数尺度是 $x^{\Re\rho}$。但这点直觉\emph{不是}本证明已用的最终零点推理。上界只有一侧，不同零点的项还会互相抵消；把每项的尺度和全体零点的结论直接连起来，正是需要避免的跳步。后面以计划好的窗口和解析准则来处理，而不是在本章就把这条积分线移到任意指定位置。


\section{把局部整除变成对偶格包含}
\subsection{先把分母真正花出去}
令 $B(x,y)=125xy$，原来的两边都是 $\Zz$。所有配对值都被 $5^3$ 整除，故可把第一边换成 $5^{-3}\Zz$：若 $x'=x/125$，新值是 $xy$，仍是整数。换格时没有改配对公式，第二格也未改变。正的整除预算对应第一格的\emph{扩大}；乘上 $125$ 虽不会破坏整性，却没用上预算。不应把秩一玩具里的 $125$ 当成真正算术中的公共元素。

在系数域 $K$ 的有限处 $w$，令 $v_w$ 是整数值的离散赋值，$\pi_w$ 是均匀化元。原格上的非零值若满足 $v_w(B_u(\phi,F))\ge q_w$，则对平移后的元素 $\pi_w^{-q_w}\phi$，
\[
 v_w\bigl(B_u(\pi_w^{-q_w}\phi,F)\bigr)
     =v_w\bigl(B_u(\phi,F)\bigr)-q_w\ge0.
\]
零值本来就整。$q_w$ 是整数，可以为负；负数的意义正是这个地点必须收缩格。若先得到的是实数尺度上的 $sg(u,w)$，源码用 $q_w=\lfloor e_w sg(u,w)\rfloor$，$e_w=v_w(p)$ 是换成整数赋值单位所需的正因子。只有既定的下界 $3$ 时，除以四个 $\pi_w$ 可能失败。若已经知道整数值赋值至少为 $3.2$，它其实至少为 $4$；采用向下取整是留有余量的统一约定，不是声称下取整给出最强的可除幂。

\subsection{这不是对所有素数做同一种运算}
$q_w$ 只在有限多个地点非零。用分式理想
\[
   I_u=\prod_w\mathfrak p_w^{-q_w},\qquad N_u=I_uN_u^0
\]
把局部平移一起实施。即使某个 $\mathfrak p_w$ 不是主理想，局部化仍能用均匀化元描述。局部证明还要把选定格上的代数配对与实际复周期配对对上。把分母在 $w$ 外清掉也不等于所有地方都已检查：$1/2$ 在 $3$ 处整，在 $2$ 处不整。

全局包含可直接验证。任取 $\phi\in N_u$ 与 $F\in R$。对每个 $w$，先将它们局部化；第一变量属于 $I_{u,w}N^0_{u,w}$，第二变量必须属于\emph{同一} $R$ 在 $w$ 的测试格。只要每处使用的值是同一 $y=B_u(\phi,F)\in K$，每处检查给出 $y\in\mathcal O_{K,w}$，取交得 $y\in\mathcal O_K$。这就证明 $B_u(N_u)\subseteq C$。若最初手里只有“各处分别存在一个 $y_w$”，从中不能直接推出存在同一个 $y$。实际源码用固定的代数配对及其指定复嵌入来辨认这些值，再经过局部格的自然性、共同格及例外地点处理。\src{Formalism/QuotientMetricDivisors/DivIntoDual.lean}{L17-L61}{对偶格内的整性}；\src{Formalism/LocalToGlobal/Divisibility/RecBridge.lean}{L128-L153}{整性在桥接数据中的位置}。

\subsection{实际链条的整除预算从哪里来}
取一个固定的 $u,w$。针对所选词，把原来的复 theta 链和局部代数链绑在同一个有理端点上：同一对象在选定嵌入处表示实际复链，在指定局部处表示选中的代数链。源码从初始输入逐步归纳：每步先用有理性取下一层的代数代表，复嵌入给出真实的下一步；同一步的交换关系则把这个代表送到局部链的下一步。归纳保持的是两个等式同时成立，不是分别从两边挑选两个彼此无关的端点。\src{Development/AbelianVarietyTheory/Proofs/SatakeParameters/LatticeValuationBound.lean}{L27-L40}{同一端点的双重代表}。这里需要追查的不只是“某条代数链有估值”，而是它就是\emph{负责这个实际周期的那条链}。达到末端后，局部列定理计入每步有效收益、附加类型项和误差。列与最终配对之间还要乘实际标量、扣掉稳定子的分母效应，并用同一有限处归一化。普通、非水平且非类型的地点，列阈值预先带的 $+vDG$ 与转成周期时的 $-vDG$ 正好抵消；因此不能把它另列成最终多扣的一项。\src{Formalism/SharpRestriction/GrowthBounds/WeilJacquetShalikaLaw.lean}{L419-L437}{实际周期估值的扣项}；\src{Formalism/SharpRestriction/GrowthBounds/WeilJacquetShalikaLaw.lean}{L504-L532}{普通地点的抵消}。把早一级列的估值原样抄成配对估值，或者每步重复扣除只出现一次的量，都可能改变答案。\src{Formalism/SharpRestriction/GrowthBounds/WeilJacquetShalikaLaw.lean}{L73-L106}{列与端点的中间式}；\src{Formalism/SharpRestriction/Divisibility/LocalQuadraticParts.lean}{L206-L221}{实际配对估值}。

关键的生成元问题来自一个容易看漏的地方。所用格不是由所有轨道一视同仁地产生；一个生成元可以只有 $S\subseteq VV$ 内的因子使用增强版，其余用基准版。固定受服务素数 $p$，记第 $j$ 步的增强下界为 $A_j$、基准下界为 $B_j$。漏掉某个本应增强的轨道，安全损失是
\[
            \sum_j \max(A_j-B_j,0).
\]
这里只计算短缺；某步的较弱目标不是借款，不把它随意用来抵销别步欠账。源码用 $p$ 的整幂补足：先逐步相加，再对该轨道的总短缺向上取整，随后把生成元上的约束延伸到所张成的格。这个取整与上节的 $q_w$ 是两件事：一个确保缺口\emph{补够}，另一个确保可花费的预算\emph{不过支}。\href{https://github.com/leanprover-community/riemann-hypothesis/issues/77}{问题 \#77} 针对的就是这一步的可读性，并未给出反例。

沿用问题 \#77 的两步算例：增强本可给出 $3.25+4=7.25$，基准只给 $1+4=5$。缺 $2.25$，额外乘 $p^2$ 仍只有 $7$；乘 $p^3$ 后达到 $8$。这里的 $3$ 是按有理素数归一化的幂数。如果固定地点上 $v_w(p)=e_w$，则在整数尺度上增加 $3e_w$，不是 $3/e_w$。不同辅助域在 $p$ 上方的地点可以有不同归一因子，换算时依赖的是同一个 $p$，不是强行令两个因子相等。这个示例只重构补偿机制，不能自行导出 $A_j,B_j$。

最后所得可用接口是：对固定 $u,w$、规定的原格元素和指定分量的测试格元素，存在相应代数值，其复像是实际配对；若该值不为零，整数赋值至少为 $q_w$。普通地点、类型地点、轨道例外、以及被排除的素数，不应靠一个普通 $p$ 的例子一并覆盖。对非零配对作估值同样不表示每个词给出的映射都非零。

\subsection{同一个元素，为什么要在同一个地方保住它的阶？}
\label{sec:same-place-order-zh}

先看一个能手算的例子。我们要找一个元素：连续乘三次 $5$ 才变成零，而且无论检查原来的整除条件，还是只检查选定的一个素因子，答案都一样。

在高斯整数 $\mathbb Z[i]$ 中，令 $\alpha=2+i$，$\beta=2-i$，则 $5=\alpha\beta$。两个因子是不同的高斯素数。取模数 $\alpha^4\beta^6$。在原商群里，变成零要同时凑足四个 $\alpha$ 和六个 $\beta$；在素理想 $(\alpha)$ 处局部化，只检查四个 $\alpha$，因为 $\beta$ 在那里可逆。每乘一次 $5$，两边各添一个因子。

\begin{center}
\begin{tabular}{@{}lccc@{}}
\toprule
元素 & 初始因子数 $(\alpha,\beta)$ & 原商群的阶 & $(\alpha)$ 处的阶\\
\midrule
$5(3-4i)=\alpha\beta^3$ & $(1,3)$ & $125$ & $125$\\
$125=\alpha^3\beta^3$ & $(3,3)$ & $125$ & $5$\\
\bottomrule
\end{tabular}
\end{center}

第一行再乘两次 $5$，只有三个 $\alpha$，还没有变成零；乘第三次，两项都恰好达标。第二行也要乘三次才能凑足六个 $\beta$，但在 $(\alpha)$ 处乘一次就够了。

这两行都从 $5$ 出发。乘 $3-4i=\beta^2$，刚好补上另一处的因子，又不改变已经有的 $\alpha$ 个数；改乘 $25$，就提前添了两个 $\alpha$，反而缩短了在选定位置的阶。正式论证里，\href{https://github.com/CheyLoveday/RH-public/blob/1c1fbd7556931388818472fb812bc3161ab12fc7/papers/type-i-torsion/lean/src/Formalism/LocalToGlobal/GlobalTheory/GL6/DiagonalSlopes/TypeISamePlaceOrder.lean\#L72-L145}{同一位置的阶数论证}用的也是这个原则：选取一个不落在指定素理想中的乘子，使关系在原格里成立，又保住前一档在该处的非零性。上面只演示这一步，并非满足 Type I 筛选条件的原始算例。

\subsection{格变大，斜率为什么恰好增加}
记 $d_u=\rk_K N_u^0=\rk_K N_u$，$l_w=\log(\operatorname{Norm}\mathfrak p_w)/[K:\mathbb Q]$ 为地点的对数权重，以下均比较同一无穷处度量下的正秩格。在整数赋值单位下，一个 $d_u$ 维的局部格乘以 $\pi_w^{-q_w}$，有符号指数为 $-d_uq_w$；例如在当地选一组基，此操作就是乘标量的对角矩阵。源码以 $2l_w$ 计这个指数在对数余体积中的贡献。用 $\mu=-\operatorname{lcv}/d_u$ 转为斜率，便得到
\[
\begin{aligned}
 \operatorname{lcv}(N_u)-\operatorname{lcv}(N_u^0)
       &=-2d_u\sum_w l_w q_w,\\
 \mu(N_u)-\mu(N_u^0)&=2\sum_w l_wq_w.
\end{aligned}
\]
这里 $d_u$ 消掉了，不可把每个格的收益再乘一次秩；筛选子格本身有共长度代价，那个代价不一定这样消掉。取整和前置删减各有明确的损失项。若直接用 $e_wsg$ 代替 $q_w$，漏掉的是小数部分带来的 $2\sum_w l_w\{e_wsg(u,w)\}$。这笔损失非负，只在实际支撑处累积。\src{Formalism/QuotientMetricDivisors/DivIntoDualShift.lean}{L18-L27}{分式理想的负号}；\src{Formalism/QuotientMetricDivisors/DivIntoDualShift.lean}{L40-L61}{斜率变化公式}。

一个在 $\Zz[i]$ 上可全算的例子：取 $\pi=2+i$，理想范数是 $5$。把 $\Zz[i]^2$ 先缩成 $\pi\Zz[i]\oplus\Zz[i]$，实余体积乘以 $5$。再将这两个坐标都乘 $\pi^{-3}$，实余体积乘以 $5^{-6}$。结果相对原格的因子为 $5^{-5}$，因此秩归一化的斜率增加 $\tfrac52\log5$；它正好是放宽所得 $3\log5$ 减去缩格成本 $\tfrac12\log5$。这也解释了为什么不能把“删掉一列”和“同倍率平移所有列”用同一个式子作账。

\section{从保留下来的方向到斜率}
\subsection{同一个映射有三本账}
整性、复矩阵和高度不是彼此独立任选的条件。取 $K=\mathbb Q(i)$，$\pi=2+i$，源格 $N=\pi^{-1}\Zz[i]$，目标 $C=\Zz[i]$，复度量都取标准度量。用恒等映射，则算子范数是 $1$，高度是零，Gram 正好为单位矩阵。源格的实余体积却缩为目标的 $1/5$，因而斜率高 $\log5$。若只凭复矩阵漂亮便说目标斜率不能低于源斜率，这里就会失败：$1/\pi$ 并不被送进 $C$。改用乘 $\pi$ 的映射，整性恢复，但用平方范数算的高度也增加 $\log5$，恰好补回。这个反例只是否定删去整性后的规则，并不针对数论中的已选族。

\subsection{迹只数总量，平方迹区分扎堆}
取四个从 $\Cc$ 到 $\Cc^4$ 的映射，每个把 $1$ 送到单位向量。第一组选 $e_1,e_2,e_3,e_4$，第二组选 $e_1,e_1,e_2,e_2$。所有映射单独看都保长度，然而第一组张成四维，第二组只张成二维。将各向量的外积相加，得到两组矩阵
\[
 Q_1=I_4,\qquad Q_2=\operatorname{diag}(2,2,0,0).
\]
两组迹都是 $4$，但 $\tr(Q_1^2)=4$、$\tr(Q_2^2)=8$。如果 $Q\ne0$ 是半正定矩阵，柯西–施瓦茨对其正特征值给出
\[
                  (\tr Q)^2\le \rk(Q)\tr(Q^2).
\]
如果目标的全维数是 $D$，并且 $\tr(Q^2)\le (1+\eta)(\tr Q)^2/D$，就有 $\rk(Q)\ge D/(1+\eta)$。在第一组 $\eta=0$，第二组最小的 $\eta=1$，两组都达到相应下界。分母 $D$ 是已知的全维数，不是未求出的像秩，否则这条准则会变得空洞。

\subsection{删掉输入之后，秩还能剩多少}
设 $V_i$ 是各个复内积源空间，$H$ 为共同目标空间，$S=\sum_i\dim V_i>0$，$D=\dim H>0$。$D=0$ 时余维为零，无需以下除法。把正权重放进矩阵：$A_i=\sqrt{c_i}B_i$，$c_i>0$。记
\[
 E_i=A_i^*A_i-I_{V_i},\qquad Q=\sum_i A_i A_i^*,
 \qquad T=\tr Q.
\]
设 $\sum_i\|E_i\|^2_{\mathrm{HS}}\le eS$，其中 $0\le e\le 1/4$；这里矩阵范数的平方是 $\tr(M^*M)$。再假设平方迹受控：\[\tr(Q^2)\le(1+\eta)T^2/D,\qquad \eta\ge0.\]删去的总维数至多 $\alpha S$，而保留的子空间为 $V_i'\subseteq V_i$。令 $W=(\sum_iA_iV_i')^\perp$，$d=\dim W$。我们要证明 $d$ 不太大。

先用 $\tr(A_iA_i^*)=\tr(A_i^*A_i)$，把目标上的迹换至各源计算。对直和上的误差块与单位矩阵施用柯西–施瓦茨：
\[
 T=S+\sum_i\tr E_i
 \ge S-\sqrt S\sqrt{\sum_i\|E_i\|^2_{\mathrm{HS}}}
 \ge S(1-\sqrt e).
\]
所以 $T>0$。接下来从两个方向估 $W$ 分到的迹 $T_W=\tr(P_WQ)$。

\textbf{上界。} 令 $R_i=(V_i')^\perp$，总维数 $k\le\alpha S$。由于 $W$ 与所有保留的像正交，$P_WA_i=P_WA_iP_{R_i}$。正交投影不增范数，因此
\[
\begin{aligned}
 T_W &=\sum_i\|P_WA_i\|^2_{\mathrm{HS}}\\
 &\le \sum_i\|A_iP_{R_i}\|^2_{\mathrm{HS}}
 =k+\sum_i\tr(P_{R_i}E_i)\\
 &\le k+\sqrt{k e S}\le S(\alpha+\sqrt{\alpha e}).
\end{aligned}
\]
这里不要求 $R_i$ 对各个 Gram 矩阵不变。被保留的方向不会给 $W$ 贡献迹，因而只有被删去的方向能供给这笔“预算”。

\textbf{下界。} 写 $Q=(T/D)I+J$，则 $\tr J=0$ 且
\[
 \|J\|^2_{\mathrm{HS}}=\tr(Q^2)-T^2/D\le\eta T^2/D.
\]
标量矩阵给 $W$ 分配 $dT/D$，实际值偏离至多 $\|P_W\|_{\mathrm{HS}}\|J\|_{\mathrm{HS}}$。故
\[
             T_W\ge (d/D)T-\sqrt{d\eta/D}\,T.
\]
令 $r=d/D$，合起来得
\[
          r-\sqrt{\eta r}
             \le \frac{\alpha+\sqrt{\alpha e}}{1-\sqrt e}.
\]
若 $\alpha\ge1/4$，直接用 $r\le1$；否则利用 $1-\sqrt e\ge1/2$、$\sqrt{\alpha e}\le\sqrt e$，再用 $2\sqrt{\eta r}\le\eta+r$，得到
\[
                        d\le (\eta+4\alpha+4\sqrt e)D.
\]
第一组四向量删去一个源线时，$e=\eta=0$、$\alpha=1/4$。较强的中间不等式直接给 $d\le1$，正好等于实际丢失的一维。第二组没有删方向，也可能出现空位，这是另一种损失。\src{Formalism/HarderNarasimhan/HarderNarasimhanProofs/FrameLinearAlgebraFirst.lean}{L96-L221}{定理的迹估计}。

\subsection{辅助行怎样投回真实周期矩阵}
在实际族中，满足迹界的矩阵可以有额外辅助行。设 $L$ 为全行空间中的保留像之和，$P$ 是投向实际行的满射。由秩零化度定理，
\[
 \dim P(L)\ge \dim L-\dim\ker P.
\]
因而 $P(L)$ 在真实目标空间内的余维，不超过 $L$ 在全行空间内的余维。形式化证明核对的并非随意截一批行，而是实际行上的矩阵确实和格映射在指定复嵌入下对应。经过两侧坐标变换，$P(L)$ 包含于实际模像的复张成空间；秩在从 $K$ 扩张到 $\Cc$ 时保持，于是可以把余维界交给格上的过滤。右侧的 $D$ 仍是全行维数，不能擅自换成真行数。若真行数为 $D'$，且辅助行的比例不超过 $\delta<1$，则 $D'\ge(1-\delta)D$，故归一化的余维上界要除以 $1-\delta$。\src{Formalism/HarderNarasimhan/HarderNarasimhanProofs/TheoremFSharp.lean}{L40-L159}{实际行投影与复秩}。

真正目标的秩也不能预先假定为正。如果它为零而 $\delta<1$，全行目标也只能为零，这会让每个 Gram 都为零。此时对角误差的总平方恰好为 $S$，与 $S>0$ 以及 $e\le1/4$ 冲突。因此目标的正秩在输入假设下得到，而非先假设所有映射都非零。

\subsection{为何选高斜率方向}
带度量的格可以有多层斜率，而不是只有一个平均值。Harder--Narasimhan 过滤按斜率从高到低列出这些层；每一个层的斜率按它的秩计重数。对本章的阅读，可以把它看作把一个单数 $\mu(N_u)$ 展开成 $d_u$ 个带顺序的数据。格的过滤本身不是对一个复矩阵做奇异值分解；它还带整数格的信息。

设 $\mu_{\max}(N_u)$ 为最高斜率，
\[
\begin{aligned}
 d_u&=\rk_KN_u,\
 s&=\frac{\sum_u d_u\bigl(\mu(N_u)-h_u\bigr)}{\sum_u d_u},\\
 U&=\max_u\bigl(\mu_{\max}(N_u)-h_u\bigr)-s\ge0.
\end{aligned}
\]
这样，从各条源斜率中扣去该词的高度 $h_u$，所得数据均值为 $s$，且任何一个不超过 $s+U$。对 $\tau\ge0$，在第 $u$ 个格只留斜率至少 $s-\tau+h_u$ 的部分 $N_u(\tau)$。若删去的比例是 $a$，则均值至多 $a(s-\tau)+(1-a)(s+U)$。因而当 $U+\tau>0$ 时，
\[
                           a\le U/(U+\tau).
\]
若 $U=0$，每条调整后的斜率都等于 $s$，没有低于阈值的方向。按魔法般的“好方向”来解释这个筛选是不够的：选择阈值的理由，是它能同时给可控的删去比例，并且在高度扣除后让所有像达到同一个目标阈值。

具体说，$B_u(N_u(\tau))$ 是 $C$ 内的模像，因为上一节已证明整性。若像非零，映射高度比较给其最低斜率至少为源的最低斜率减 $h_u$，因而至少为 $s-\tau$。过滤的极大性把它包含在 $C$ 同一阈值的部分 $C(\tau)$。若像为零，包含自然成立，无需对零映射取对数。这样所有保留像的和都在同一个 $C(\tau)$ 里，上一节的余维上界也就约束了 $C(\tau)$ 的余维。高秩绝不意味着 $B_u$ 各自单射，也并不要求每个词贡献非零。\src{Formalism/HarderNarasimhan/HarderNarasimhanProofs/TheoremFSharp.lean}{L277-L337}{实际过滤与秩界}。

\subsection{沿阈值积分，得到平均斜率界}
写目标格斜率为 $b_1,\ldots,b_{D'}$，按秩计重数。取 $H_0>0$，覆盖所有 $s-b_k$ 的正部。逐项积分指示函数得到
\[
 \frac{1}{D'}\sum_k(s-b_k)_+
 =\int_0^{H_0}\frac{\#\{k:b_k<s-\tau\}}{D'}\,d\tau.
\]
例如两个目标斜率是 $s-2$ 和 $s+1$，则被积分比例从 $0$ 到 $2$ 是 $1/2$，之后为零，积分为 $1$。真实均值低于 $s$ 只是 $1/2$，因为正部不计高于阈值的抵消；因此用积分做上界方向是正确的。

代入余维上界并按 $\tau$ 积分，得到
\[
 {\cal E}_{\mathrm{fr}}
 =\frac{4U\log(1+H_0/U)+(\eta+4\sqrt e)H_0}{1-\delta}.
\]
$U=0$ 时第一项用连续延拓取零。因为 $s-\mu(C)\le D'^{-1}\sum_k(s-b_k)_+$，故在该接口下 $s\le\mu(C)+{\cal E}_{\mathrm{fr}}$。底层源码确有这条强式；实际组合层选用的是加了 $\log|\mathcal W|$ 的弱式，$\mathcal W$ 为词族。这里记录这一区别是为了忠实追踪主链，并不据此声称已经得到新的素数误差或改进指数。\src{Formalism/HarderNarasimhan/HarderNarasimhanProofs/TheoremFSharp.lean}{L338-L400}{积分和强弱两式}。

\section{真实周期矩阵的平方迹从哪里来}
前一节的秩引理只需要一些内积和矩阵。它不会替你证明，真正从 theta 提升链来的矩阵就具有那么小的误差。这里把源码中的两路输入拆开看：一路认出正在估的矩阵，另一路提供相关性的大小。只有到两路接上，形式化的 Frame 结果才能用于真实格。

\subsection{第一路：辨认矩阵，而不是估范数}
为免与上一节实际源秩 $d_u$ 混淆，记未压缩块为 $P_u:\Cc^{b_u}\to\Cc^D$，压缩后的块为 $P_uV_u$，其中 $V_u$ 是从 $d_u$ 维源空间到 $b_u$ 维空间的等距嵌入。$D$ 包含辅助行；真实行数记作 $D'\le D$，另一个出现在点计数和归一化中的量叫 $D_\star$。这三个数不预先相等。

BesselId 在这里不是一条从非负性突然推出的范数小结论。它确认两种方式写出来的配对是同一个。其一从数域配对 $B_k(u)$ 出发，在选定无穷处嵌入；其二从已有的周期配对出发，把输入放入 Hermite 空间，再用剪前块计算。源码中的复配对可以写成
\[
 B_u(\phi,F)=a_u\,
     \operatorname{bes}_u(U_u\phi,VF),
 \qquad a_u\ne0.
\]
这里 $U_u$ 随 $u$ 变化，$V$ 是整个族共用的目标侧变换；两者在此接口中是酉变换。前式只识别配对，非零系数 $a_u$ 不是下节的正归一化权 $\gamma_u$。

证明中的辨认不是比较两个空洞的名字。它在同一 $(u,t)$ 上使用实际运输的周期列，与选定 theta 像的 Hermite 配对作比较。$t$ 记检验端的一个索引。再经各分量的代表和拼接，等式才得到对规定全部输入的版本；数域配对与复周期的识别是另外一项，最终并入同一接口。真实行上的剪后块这时才可写作
\[
                  Z\,\sigma([B_k(u)])\,Y_u ,
\]
其中 $Z$ 是公共目标坐标变换，$Y_u$ 是各源端变换并吸收 $a_u^{-1}$。$\sigma$ 取指定复嵌入，源码也处理经复共轭的坐标表达。公共 $Z$ 是必要的：假如每个字都任意转动目标坐标，逐块的 Gram 不变，叠在一起的张成空间却可能变小。两组四向量例子恰好警示了这点。\src{Formalism/LocalToGlobal/IndefiniteReal/IsometricPeriodsOfHypotheses.lean}{L97-L106}{运输后实际配对的等式}；\src{Formalism/LocalToGlobal/Divisibility/RecBridge.lean}{L327-L332}{数域配对与 BesselId 的合成}；\src{Formalism/SharpRestriction/GlobalSchurBounds/TopDensityNGeSixteen.lean}{L1283-L1315}{实际行与公共坐标}。

\subsection{第二路：哪些算术界需要另给}
取来源规定的正权 $\gamma_u=1/c_0(u)$，在压缩前定义交叉块
\[
            K^f_{uv}=\sqrt{\gamma_u\gamma_v}\,P_u^*P_v.
\]
LinkS 给的是相互区分的三条界：
\[
\begin{aligned}
 \|K^f_{uu}-I\|_{\mathrm{HS}}^2&\le e_0 b_u,\\
 \|K^f_{uv}\|_{\mathrm{HS}}^2
       &\le\frac{(1+\eta_0)b_ub_v}{D_\star}\quad(u\ne v),\\
 D&\le(1+\eta_1)D_\star.
\end{aligned}
\]
第一条说各字分别不坏，第二条才说字与字放在一起也不太坏，第三条把计数尺度与矩阵行数对齐。辅助行比例的 $\delta$ 由 LinkIRR 的行数控制，不是这里的源端删减。要另有维数下限与可控的源方向删减，否则坏掉的行列比例会把这些精度吃掉。把 $D_\star$ 一开始当作真实行数，会悄悄移除第三条的证明任务。\src{Formalism/SharpRestriction/GlobalSchurBounds/TopDensityNGeSixteen.lean}{L817-L852}{对角、非对角和行数输入}。

非对角那一条的上一层并不是重复做一次 Cauchy--Schwarz。源码经核展开、收敛和消失性，将块的平方范数对应到特定的主核；然后调用针对已选参数的非对角核界。对角和点计数经另外的接口组合进来。这一节已追到它们因何有权谈论同一组矩阵；\emph{没有独立重证主核的定量估计及上游选族}。若写整部传统证明，那里应写成能从原始数据构造核并推出界的章节，而不只是把 LinkS 整句复制进来。\src{Formalism/LocalToGlobal/SiegelLatticeTheory/FlagStabilizerTensors/PackageTie.lean}{L225-L235}{块平方范数的核公式}；\src{Formalism/LocalToGlobal/SiegelLatticeTheory/FlagStabilizerTensors/LinkSFinalP.lean}{L98-L114}{非对角的定量输入}；\src{Formalism/LocalToGlobal/SiegelLatticeTheory/FlagStabilizerTensors/LinkSPos.lean}{L149-L161}{对角界与计数的接口}。


\subsection{异字核界再向上追了一层}
沿实际选中族往上走，实际入口是 LinkS 的「零项已经处理」分支。异字核估计在这个入口由行估计送到核上，再调用专门对不同词成立的定理。这一点要紧：仓库里相近的通用包装不一定是实际执行路径。这里只称“异字”的是两个不同的词 $u\ne v$，不是说所有坐标向量天然不相关。\src{Development/AutomorphicArithmetic/ThetaCorrespondence/ThetaPairLifts/ThirdPlaceSeries/LinkAltAltBis.lean}{L39-L47}{选定主链的调用}；\src{Development/AbelianVarietyTheory/Proofs/FlagMeasures/OffK0Closed.lean}{L673-L683}{异字输入的实际构造}。

异字所做的一件实质性工作，是让一个退化核部分消失。另一方面，恒等胞腔的退化条件对所有配对都另行成立，不能把它和 $u\ne v$ 的结果混为一谈。二者移除的是不同项。“异字”本身并不从逻辑上推出“小”字；剩下的逐类大小要靠定量估计。\src{Formalism/LocalToGlobal/SiegelLatticeTheory/FlagStabilizerTensors/PackageTie.lean}{L225-L235}{消失条件怎样进入核平方公式}。

具体估的是带权内积的绝对值之和，而不是要求每条向量的范数都小。源码先把贡献切成几类。亏秩至少二的部分，以及泛型亏秩一中不属于标量主项的部分，交给绝对上界；在那里先用三类积分或格点估计，控制各种退化情形。经参数比较，形如 $\exp(-kf\log k/2)$ 的项乘上前因子后，仍可压到 $\exp(-kf\log k/3)$ 一类的界。此步不是对每条向量单独算范数，而是先控制类内总量，再按并集分给各个误差项。\src{Formalism/HermitianCurvature/StatementPB.lean}{L400-L408}{绝对界的输入形状}；\src{Formalism/HermitianCurvature/HermitianCurvatureProofs/GrowthOfBasic.lean}{L1191-L1228}{格点估计的组合与指数预算}。

主项并没有全部消失。泛型亏秩一中的标量部分有一个保留的行主项；非泛部分则用另一个误差约束。到最后一层才利用归一化、主项偏差和其它小量，把这几类加成带 $1+b^{-6}$ 的目标界。这个 $1$ 不是记账疏漏：目标本来也不是让异字的每一对内积都严格为零，而是控制平方范数的总量与维数尺度之间的比例。把所有异字项一律设成零，反而不是这份证明。\src{Formalism/HermitianCurvature/HermitianCurvatureProofs/GrowthOfBasic.lean}{L464-L550}{主项保留和各类合并}。

现在可以明确说，剩余难点不在最后的指数代数，也不在上节的平方迹恒等式。所选族先各自取得统一的局部计数常数和大间隙振幅界，再与小胞腔体积、项编码的单射及固定格的条件合并；这是代码里把抽象上界用到实际族的关口。\src{Development/AbelianVarietyTheory/Proofs/FlagMeasures/LoopRowBoundsAlongFamilyLargeGaps.lean}{L25-L43}{所选族的统一估计}。需要独立写出的是三类上界为何对实际选中局部对象成立，尤其是大间隙下的 Bessel／Whittaker 控制、格点和与余体积如何交叉使用。若这些部分只能引用 Lean 声明，中段教材化仍能成立，但不能把整部书称作可脱离源码逐条重证的纸笔证明。

\subsection{三条算术界合成一个矩阵误差}
令 $A_u=\sqrt{\gamma_u}P_uV_u$，$Q=\sum_u A_uA_u^*$，$T=\tr Q$。循环移迹给出
\[
                \tr(Q^2)
                =\sum_{u,v}\|V_u^*K^f_{uv}V_v\|^2_{\mathrm{HS}}.
\]
推导只需将 $Q^2$ 展开，并将 $A_u^*A_v$ 与它的伴随移到迹号内。等距嵌入对应的压缩不放大每项。设 $R_s=\sum_u b_u$。用三角不等式，剪前对角项总量不超过 $(1+\sqrt{e_0})^2R_s$；对非对角项求和则不超过 $(1+\eta_0)R_s^2/D_\star$。两类贡献尺度不同，不能从第一项小推出第二项也小。用 $D_\star/\mathrm{sg}\le R_s$，对角总量才能按二次尺度一起入账：
\[
 \tr(Q^2)\le
 \Bigl(1+\eta_0+\mathrm{sg}(1+\sqrt{e_0})^2\Bigr)
              \frac{R_s^2}{D_\star}.
\]
压缩后的对角误差满足
\[
 \sum_u\|V_u^*(K^f_{uu}-I)V_u\|^2_{\mathrm{HS}}\le e_0 R_s.
\]
因为 $V_u^*V_u=I$，压缩前与压缩后对恒等的偏差能在这里相接。若实际源秩总和至少 $(1-\mathrm{de})R_s$，则按实际源秩归一的误差最多为 $e_0/(1-\mathrm{de})$。这一步把误差转进了前节的 $e$，而不是沿用剪前平均值。

再设压缩前后损失的源维数不超过 $\mathrm{de}\,R_s$。剪前总迹至少 $(1-\sqrt{e_0})R_s$。用前节的投影方法减去被删的源方向，它们带走的迹不超过 $(\mathrm{de}+\sqrt{\mathrm{de}\,e_0})R_s$。于是
\[
   T\ge(1-\sqrt{e_0}-\mathrm{de}
                   -\sqrt{\mathrm{de}\,e_0})R_s.
\]
令 $\beta=1-\sqrt{e_0}-\mathrm{de}-\sqrt{\mathrm{de}\,e_0}$。在 $\beta>0$ 下，代入行数比较得到可直接复算的上界
\[
 \frac{D\tr(Q^2)}{T^2}
 \le\frac{(1+\eta_1)
       (1+\eta_0+\mathrm{sg}(1+\sqrt{e_0})^2)}{\beta^2}.
\]
再由 $\eta_0,\eta_1\le\mathrm{sg}/2$、$\sqrt{e_0}\le\mathrm{sg}/2$ 得到
\[
 \tr(Q^2)-T^2/D\le\eta T^2/D,\qquad
 \eta=5x^{-CH}+3\sqrt{\mathrm{de}} .
\]
这里 $CH$ 是该段的参数，$\mathrm{sg}=x^{-CH}$，不能把出现在常数中的 $\mathrm{de}$ 当作某一次过滤丢掉的比例 $\alpha$。前者是源矩阵建模时的压缩损失，后者是已经得到实际格以后，在斜率阈值上再删掉的方向。以上推导解释了为何真正要查的是非对角核与计数，而不是试图在已成立的有限维迹不等式里找出凭空的数论。\src{Formalism/SizeWindowBrackets/ChainSlopeWindows/TraceConjTransposeMulSelf.lean}{L549-L625}{迹和压缩所用的界}；\src{Formalism/SharpRestriction/GlobalSchurBounds/TopDensityNGeSixteen.lean}{L602-L646}{实际小量与最终参数}。

\section{两次相消，留下真正待估的量}
\subsection{配对的高度不等于你想证明的素数界}
在固定族点，有限处一侧给平移后源格的斜率下界，形状为
\[
   \mu(N_u)\ge\mu(M_u)+G_a+F_o-E_e(u)-L_o(u).
\]
$M_u$ 是基准格，$G_a$ 来自局部有效收益的总账，$F_o$ 是同一族点上换格斜率下界与高度估计共用的修正量，而 $E_e,L_o$ 是另外处理的成本。这行按同一归一化将局部重数汇入斜率。这个把各地点按有理素数合并的步骤，必须按对应地点的权重进行，不能按一个 $p$ 上方的地点数再乘一次。也不能把 $G_a$ 定义成目标主项，从而不经过比较。

除法及无穷处估计给出\emph{同一个} $h_u$ 的上界：
\[
       \mu(C)+h_u-F_o\le \mathrm{MAIN}_u+E_{\mathrm{div}}.
\]
移项可见 $h_u\le\mathrm{MAIN}_u+E_{\mathrm{div}}-\mu(C)+F_o$。两式相减，有限处一侧的 $+F_o$ 与高度里的 $+F_o$ 相消，得到
\[
\begin{aligned}
 \mu(N_u)-h_u\ge G_a+\mu(C)+\mu(M_u)
  -\mathrm{MAIN}_u\ &{}-E_e(u)\\
  &{}-L_o(u)-E_{\mathrm{div}}.
\end{aligned}
\]
若把高度式左侧写成 $\mu(C)+h_u+F_o$，就不会得到同一结论。若 Frame 中实际用的不是同一个 $h_u$，那么按上式计算的加权均值也不是 Frame 要比较的那个数。\src{Formalism/SharpRestriction/GlobalSchurBounds/QuadraticLogDerivCongr.lean}{L580-L591}{高度估计中的公共修正}；\src{Formalism/SizeWindowBrackets/ChainSlopeWindows/LevelOfCongruence.lean}{L50-L66}{按实际秩逐词平均}。

\subsection{逐词的主项如何变成公共主项}
前式的 $\mathrm{MAIN}_u$ 不是常数。令 $A=\sum_t a_t\ge0$，其中含有随词改变的一项 $-A\log T_u$。计数输入给出同一族点上的共同下界 $\log T_u\ge\Lambda_5$，所以在估高度上界时，可以用 $-A\Lambda_5$ 代替原项。$A\ge0$ 对方向重要。另一个输入估计 $E_e(u)$，其主项为 $mA\ell_g$，并把差放入误差。于是高度与斜率预算的公共部分汇成
\[
 S_{\mathrm{Div}}-A\Lambda_5+mA\ell_g.
\]
这里 $\ell_g$ 是源码 \texttt{lg}，$S_{\mathrm{Div}}$ 是 \texttt{SDiv}。对应恒等式是
\[
\begin{aligned}
 S_{\mathrm{Div}}-A\Lambda_5+mA\ell_g
      &=\mathrm{Main}+A E_{\mathrm{count}},\\
 E_{\mathrm{count}}&=c_{\mathrm{count}}
      \bigl(\mathrm{La}+m^2\mathrm{La}/k\bigr).
\end{aligned}
\]
一眼看去，这三项往往都像“调参”。实际职责不同：$S_{\mathrm{Div}}$ 随除法高度进入，$\Lambda_5$ 由共同的水平计数下界进入，$mA\ell_g$ 来自指数成本的主项。公共修正相消之后，再把这三项用具体代数恒等式整理成 \texttt{Main}，不是预先拿素数偏差的上界当输入。\src{Formalism/SizeWindowBrackets/ChainSlopeWindows/LevelOfCongruence.lean}{L17-L27}{水平计数和预算定义}；\src{Formalism/SizeWindowBrackets/ChainSlopeWindows/LevelOfCongruence.lean}{L105-L111}{主项恒等式}。

其余的基准斜率、指数界与 $L_o(u)$ 都要付进去。按实际源秩 $d_u$ 平均后，上述推理得到
\[
       s\ge G_a+\mu(C)-\mathrm{Main}-\mathcal E_{\mathrm{arith}}.
\]
这里 $\mathcal E_{\mathrm{arith}}$ 是为了叙述而合并的数，不是声称源码用同名符号。它包括水平计数余项、除法和源格的估计，以及按 $d_u$ 平均的损失。若把此处改为对词作等权平均，得到的通常就不是同一个 $s$。例如两个词的秩为 $1$ 和 $9$，调整后各为 $10$ 和 $0$，等权均值是 $5$，按秩均值却是 $1$。

\subsection{第二次消去是真的，不是循环}
Frame 上界与刚得到的下界可并列：
\[
\begin{aligned}
  G_a+\mu(C)-\mathrm{Main}-\mathcal E_{\mathrm{arith}}
      &\le s,\\
  s &\le \mu(C)+\log|\mathcal W|+\mathcal E_{\mathrm{fr}} .
\end{aligned}
\]
消去的是同一个 $s$，随后消去这同一个对偶格的 $\mu(C)$。结果为
\[
 G_a-\mathrm{Main}\le\log|\mathcal W|
      +\mathcal E_{\mathrm{arith}}+\mathcal E_{\mathrm{fr}}.
\]
这既说明为什么不需要单独精算 $\mu(C)$，也说明不能跳过确认“到底是哪一个 $C$”。下界使用局部整除、换格和除法高度；上界使用同一映射的整性、矩阵与斜率过滤，来源与结论并不自相循环。进一步还要单独控制 \texttt{ASK}$-G_a$ 及轨道损失 \texttt{IO}，才得到应用于窗口平均的那条估计。即使 \texttt{ASK}$-G_a$ 在某个族点碰巧为零，也不是整个证明可把它删掉的理由。\src{Formalism/SizeWindowBrackets/ChainSlopeWindows/LevelOfCongruence.lean}{L113-L197}{消项、Frame 误差和点态估计}。

\subsection{窗口是输出接口，不是任意点的结论}
本章结束在算术上界的交接点。对族点的估计、对窗口的平均、以及对大尺度解析对象的结论，是三个不同层次。后续解析环节在指定参数与薄尺度条件下，才把上界转为含 $x^{1/2+O(\varepsilon)}$ 的估计，并处理剩余的特征项。这里不能把存在可用窗口写成任意点结论，不能把单侧直接写成绝对值界，也不能把依赖 $\varepsilon$ 的阈值假装与 $\varepsilon$ 无关。\src{Development/AutomorphicArithmetic/DedekindLFunctions/LandauTheorems/ChCore.lean}{L343-L351}{中间解析命题的输入和量词}。再往下的一个明确接口是 \texttt{WindowHyp}。对固定 $\theta,\delta$ 和字符族，它要求对每个 $\eta>0$ 与每个起点下限 $Y_0$，都能选 $Y\ge Y_0$ 和一个非负上盖 $R$，令混合偏差在 $[Y,Y^{1+\delta}]$ 上逐点不超过 $R$，而且
\[
       \int_Y^{Y^{1+\delta}}R(y)y^{-\theta-1}\,dy\le Y^\eta .
\]
可积性是单列的条件。这是“任意靠后有够长的可用窗口”的准确量词，不是说任意这样的窗口都已受控。\src{Development/AutomorphicArithmetic/DedekindLFunctions/CharFamily.lean}{L38-L42}{窗口假设与加权积分}。

书稿若要覆盖 RH 的完整论证，需把这里再接到精确的非平凡零点判据，而不是仅把 $S(x)$ 放回结尾。

\section{500 页能装下什么}
本章实际完成的，是把中段依赖拆成读者可以复算的局部：受服务素数的补偿和分母方向、全局整对偶格、有限维秩不等式、阈值积分、两次抵消。读者无需读数千条源码辅助声明，就能看出换掉 $C$、$h_u$、或者 $B_u$ 的任何一个会在哪一行失败。这说明压缩不必损失审计性。

但这段能写成试读稿，不等于整部证明已经可以独立压成 500 页。真正的计价单位应是“尚未重证的输入”，而不是 Lean 的行数。以下是以数论研究生为读者、允许明引标准定理、只写经典 RH 的试编预算：
\begin{center}
\begin{tabularx}{\linewidth}{@{}Xr@{}}
\toprule
职责 & 页数预算\\\midrule
解析判据、窗口、辅助域与标准工具 & 120\\
局部提升、整除、平移与 Frame，两侧估计 & 230\\
Galois、Sha、余体积等算术输入 & 85\\
全局装配、量词、依赖和源码索引 & 65\\\midrule
合计 & 500\\\bottomrule
\end{tabularx}
\end{center}
这不是已经实现的页数承诺。最应先压测的是两个大输入：真实局部 theta 链如何逐例外地产生可用下界，以及真实周期族的非对角 Gram 控制如何得到足够小的量。前者若只列结论，读者看不到素数收益从哪里来；后者若只给个等距例子，读者看不到方向为什么不扎堆。就是否可读而言，能否写清这两处比减少记号或合并小引理更重要。

\subsection*{书稿的验收办法}
请有代数数论基础的读者不看 Lean，先读两条支路和合流处，能自己回答以下问题：
\begin{enumerate}
 \item 为什么补偿向上取整、平移向下取整？分别在哪种单位上算？
 \item 为什么大迹不够，平方迹还需要完整目标维数？
 \item 哪三条事实保证过滤后的像进的是实际的 $C(\tau)$？
 \item 如果把第二次相消中的 $C$ 改成同秩异度量格，哪行先错？
 \item 哪些新引理靠书稿的证明，哪些依赖源码或外部已有定理？
\end{enumerate}
倘若答不上最后一问，这只能算一份导览，不能称独立证明。若前四问可答而第五问有明确的依赖表，则已达到“读懂路线并追查关键环节”的标准。经典 ζ 情形先独立成书；Dirichlet GRH 可以做后续卷，不抢占这 500 页的预算。

\subsection*{阅读对照}
未另标来源的演示与算例由本次审读组织，源代码用于核对精确接口和依赖。对照材料包括 \href{https://www.math.uwaterloo.ca/~mrubinst/alpoge_claude_grh/zeta-paper.pdf}{Rubinstein 的 ζ 情形试稿}（尤其 III.5--III.8）、\href{https://tdupu.github.io/riemann-public/riemann-book.pdf}{Dupuy 的讲义与输入表}（表 3.3）以及 \href{https://leanprover-community.github.io/blog/posts/riemann-hypothesis/}{形式化项目联合说明}。所有源码链接均指向本章所审的固定版本，不据草稿尚未展开的输入断言证明不存在。

\section{复算清单与依赖表}
\subsection{三道短题可分辨哪里出了偏差}
\textbf{题一：三次取整是否同一件事？}
固定一个受服务 $p$。轨道甲两步各欠 $0.2$，轨道乙两步各欠 $0.6$。先在各轨道内部按步求和再上取整，分别需要补 $1$ 与 $2$ 重。若在每步上取整，得到的是 $2$ 与 $2$ 重，会付出更多；若把两个轨道合成一个再取整，只得到 $2$ 重，不能用这个总数同时满足两项独立的轨道约束。最后换算到某个 $w$ 的总下界后才由 $q_w$ 向下取整。注意最后一步甚至可能得到负数，表示收缩原格，不再是在补哪条局部轨道。

\textbf{题二：共用坐标变换哪里真正有用？}
四个源都为一维，目标在数域上有四维。假如四个算术像都恰好落在第一根坐标轴，把每个字的目标端各自转到不同坐标轴，就会让辅助矩阵看似张成四维，而真实像的和只有一维。若整个族共用同一个可逆 $Z$，则
\[
      Z\Bigl(\sum_u \operatorname{im}B_u\Bigr)
      =\sum_u\operatorname{im}(ZB_u),
\]
总维数保持。源侧可以每字不同，因为可逆源换基不改该字的像。实际证明把这一点进一步用于过滤后的相应子空间，然后用真行投影将辅助行扣掉。故共享目标侧换基不是一个格式要求。

\textbf{题三：为何不能把对偶格换成一个同秩格？}
在前面的 $\Zz[i]$ 秩一例子里，$C$ 和 $\pi^{-1}C$ 秩相同，但默认实度量下斜率相差 $\log5$。如果一条不等式给 $s\ge G_a+\mu(C)-M-E_1$，另一条只给
$s\le\mu(\pi^{-1}C)+E_2$，就只能推出
$G_a\le M+E_1+E_2+\log5$，而不能把多出的 $\log5$ 消去。这个例子不说源码真的换错；它告诉审读者该在拼接处比较格的身份和度量，而非只比较秩。

\subsection{四个尚须用传统数学语言展开的接口}
\textbf{局部算术。} 列下界的基准分支和扭转分支应分开书写，说明由哪个等价关系把局部代数列落回到实际配对。正文已将生成元补偿与全局理想贴合分开，但没有从局部 theta 核逐边重算各类素数的端点估值。用同一有理 $p$ 比较辅助域处与系数域处时，还须明列各自赋值单位。\src{Formalism/SharpRestriction/GrowthBounds/WeilJacquetShalikaLaw.lean}{L19-L50}{列界的对象、分支和例外}。

\textbf{周期的非对角界。} 在矩阵层，核的恒等式只告诉读者应该估谁，不能自己给出小量。将异字展开为各类项，解释为什么相同路径可进主项、真正异字需要消失或衰减，以及例外计数如何叠入余项，是更长书稿应做的工作。本章用 LinkS 给定的数值界推到中心平方迹，并未把核的算术生成与分析估计全盘重证。

\textbf{主项高度。} 消项所需的代数恒等式很短，但它从 $S_{\mathrm{Div}}$ 中提取的是已证明的高度主项。完整篇幅不能仅说“恰好抵消”，应写出除法高度、水平点计数及指数成本各自的来源，给同一个族点的参数对照，再算出该恒等式。局部与高度两条界都调用了 $F_o$ 的同一个定义；来源相同不意味着把一条界当作另一条重新使用。

\textbf{最终量词。} 即使中段可以完全独立阅读，仍要从几种窗口和它们的成本比较，到所需解析估计以及目标零点结论逐级交接。最清楚的写法是先写固定精度，再写阈值依赖、在何处取平均、哪些大尺度条件必须满足；最后才使用“任意精度”。\src{Development/AutomorphicArithmetic/PrimitiveCosets/CoreGm.lean}{L12-L22}{后续解析比较如何陈述条件}。

\subsection{源码审读时的身份索引}
从实际端点到最终不等式，有几处即使名称相同也不能只靠拼写判断：
\begin{itemize}
\item $p$ 与 $w$：有理素数及其上方的地点。支持集、付出的幂、以及整数赋值中的贡献要分开。
\item $m,b_u,d_u$：迭代步数、剪前复源维数、实际源格秩。Frame 平均用第三个；最早块界用第二个；局部合计常用第一个。
\item $D,D',D_\star$：全行、真行、算术归一化尺度。三个量通过已证明的比较交接。
\item $\alpha,\mathrm{de},\delta$：实际过滤删掉的源秩比例、建模时剪掉的源秩比例、目标辅助行比例。它们不是一次筛选的三个名字。
\item $B_k,B_u,P_uV_u$：分别是数域格上的映射、实际复配对、全行矩阵。它们的身份通过复嵌入、坐标变换和真行投影联系，不是一个声明本身。
\item $G_a,\mathrm{ASK},\mathrm{Main}$：局部有效收益、待比较的量与主项。$\mathrm{ASK}-G_a$ 的控制在斜率合流后仍有任务。
\end{itemize}
一个可持续的人工审读流程，应让每个新引理能回到这张表上的至少一个具体接口。若某次压缩删了定义但没有留下足够的信息辨认输入，省下的页面要在返工时重新付出。

\clearpage
\section{哪一项对所有输入都为真？}

\noindent\textbf{先给结论。} 证明中确实有一个不给输入设限的几何条件。它被最终证明引用，但负责核验定理的程序并不是这段代码。要分清三件事：哪一项为真、其余条件从哪来、最后核验的是什么命题。

\subsection{恒真的只是几何一栏}

Jenny Lorraine Nielsen 在批评稿中指出，原代码里有这样一行：
\begin{quote}
\small\verb|def anyGeometry : GeoP := fun _ _ _ => True|
\end{quote}
它有三个参数，三个下划线表示一概不用。无论代入哪组对象，得到的命题都是 $\mathrm{True}$。同一文件的 \texttt{geoOf} 则不同：它要求一个正整数和满足指定条件的几何见证，同时要求 \texttt{HasTrueSurf}。两种条件可以填进同一个位置，但用前者不会给出后者的见证。
\src{Development/AbelianVarietyTheory/Definitions/AutomorphicMembership/OverIx.lean}{L16-L19}{查看定义}；
\href{https://philpeople.org/profiles/jenny-lorraine-nielsen/publications?order=viewings}{查看批评稿摘要}。

代码把这个条件包在 \texttt{GeoCl} 里。它的基本构造只要求先给出底层条件的证明。底层若是 $\mathrm{True}$，基本构造对任意类型相符的输入都能填上。这里确实没有筛选，但筛的是一项几何条件，不是零点。就像把一道题的一栏标为「无需填写」，不能据此把其余各栏也算作答完。

\subsection{最终证明用到了它}

沿最终定理往回追，确实能走到这项恒真条件。原码在证明「存在一组满足后续条件的数据」时，把它传给了组装这些数据的三处调用。两份依赖审计也找到了这条引用链。Nielsen 点名的 \texttt{plugAll} 不在所审计版本的最终链上；这不改变恒真条件确实被引用的事实。
\src{Development/AutomorphicArithmetic/ThetaCorrespondence/ThetaPairLifts/ThirdPlaceSeries/SiegelEisensteinSecondTermHypotheses.lean}{L44-L73}{查看调用}；
\href{https://github.com/harrytong/rh-verification/blob/main/RH_critics_dependency_findings.md}{查看依赖审计}。

\newpage
\subsection{真正要追的是其余前提}

同一个构造还接收另外几项有内容的前提。例如，共体积是一个格的基本胞所占的体积；在二维里，就是两条基向量张成的平行四边形面积。$h_X$ 给出的界，控制每个合格数据族在尺度足够大时的归一化对数共体积。代码把这条界用在选出的各个分量上。另两项前提分工不同：$h_M$ 保证合格数据存在实处和有限处的匹配；$h_A$ 控制匹配后两边 $\theta$ 积分的比例关系等，并不要求它们直接相等。构造检验向量和取值网格都用到了这两项。

这些数据组装好以后，代码再用一条引理排除狄利克雷 $L$ 函数在临界线右侧的零点（极点另行排除），最后推出普通黎曼猜想。读这一段时，应追问共体积界从哪里来、局部匹配怎样进入整体估计，以及哪条引理把估计变成无零点结论。不能把「几何一栏通过」写成「已经找到了实际的几何见证」。

\subsection{怎样排除「核验器照单全收」}

核验者只用 Mathlib 另写了待证命题，再与这份代码推出的定理比较。两边匹配的是原来的黎曼猜想：$\zeta$ 的每个非平凡零点实部等于 $1/2$，不是 $\mathrm{True}$。独立检查器复核了静态导出的证明项，不运行项目自带的证明生成代码。五条主要定理使用的公理都是 Lean 的三条标准公理。
\href{https://leanprover-community.github.io/blog/posts/riemann-hypothesis/}{Lean 与 Mathlib 联合声明}；
\href{https://math.mit.edu/~drew/grh/}{MIT 核验记录}。

还有一个对照实验：Carlos Toledo 把导出中普通黎曼猜想定理的目标类型改成 $\mathrm{False}$，保留原证明项不变。两套检查器 \texttt{nanoda} 和 \texttt{con-leche} 都拒绝了这份改过的输入。这检验的是检查器能否拒绝错误，不代替前面那些数学步骤的解释。
\href{https://carlostoledo.co/reports/alpoge-grh.html\#changes}{11 月 22 日实验记录}。

\end{document}
